\section{Through Advent}

\subsection{The First Sunday}

The entrance chant takes its words from a psalmist who has enemies and asks
God to keep him from shame. “To thee, O Lord, have I lifted up my soul” (Ps.
24:1). This is the Introit, the chant at the entrance. The station named in
the Missal is St Mary Major. A station is a church in Rome appointed for the
celebration of Mass on a particular day. In today's Collect, the prayer the
priest says before the readings, he asks the Lord to stir up his power and
come to our rescue. The danger the prayer speaks of comes from our own sins.

St Gregory the Great looks closely at the different ways people meet the
terror Jesus describes (Homily 1 on the Gospels, §§1, 3). Whoever has put his
hope in the present world sees it giving way beneath him. For the disciple it
is the one he loves who draws near. To lift the head, Gregory says, is to
raise the mind toward the joy of heaven. He dwells on the short parable Jesus
tells after speaking about the sea and the stars. When trees put forth their
fruit, you know summer is near. Gregory hears in that ordinary knowledge,
which anyone who has watched a tree can have, the promise of a day on which
the clouds of sorrow will be gone.

St Paul speaks at the hour when night is nearly over and it is time to get up
(Rom. 13:11--14). He tells the Romans to lay aside works of darkness,
including quarrelling and envy. The last words read at Mass tell us what to
put on for the coming day, the Lord Jesus Christ himself. If there is a duty
you have been avoiding because you would rather have someone else do it, begin
it today. A person to whom you owe an apology may be within reach before
nightfall.

\subsection{The Second Sunday}

The station is the Roman basilica of Holy Cross in Jerusalem. John is in
prison. He has heard of Christ's works and sends two disciples with a
question. “Art thou he that art to come, or look we for another?” (Matt.
11:3). Jesus tells them to report what they are hearing and seeing. Blind
people see, lame people walk, and poor people hear the good news. Anyone who
has heard the prophet Isaiah speak of eyes opening and the lame leaping, and
of good news brought to the poor, has some idea of what such deeds mean (Isa.
35:5--6, 61:1). St Paul says in the Epistle that the Scriptures were written
for our learning, so that we might have hope (Rom. 15:4).

The last thing Jesus says to John's messengers is “blessed is he that shall
not be scandalized in me” (Matt. 11:6). In this promise, to be scandalized is
to stumble over Christ himself. Gregory pictures the Baptist at the threshold
of his own death, asking whether the one whose coming he announced among the
living will also come among the dead (Homily 6 on the Gospels, §1). He takes
the warning to refer to the Passion. The people who were amazed to see Christ
give life would see the life-giver put to death. Gregory sees the love for
which Christ is worshipped at its deepest when the life-giver goes willingly
to his death. “For the more humbling God hath undergone for man's sake, the
more worthy is He that man should worship Him.”

After the messengers leave, Jesus asks whether the crowd went into the desert
to see a reed shaken by the wind (Matt. 11:7). Gregory thinks of the way
praise and blame can blow a person in opposite directions (Homily 6, §2). John
continues to serve the one he has announced even when his preaching has landed
him in prison. The Collect asks God to stir up our hearts so that we can
prepare the way for his Son. Ask God to stir up your heart. Do the good you
would be tempted to leave undone because no one will thank you.

\subsection{The Third Sunday}

“The Lord is nigh” (Phil. 4:5). Paul can write this from prison, where his own
bonds have made known the reason he is there and given others courage to speak
(Phil. 1:12--14). The station is St Peter's. The Introit begins with his
command to rejoice, \emph{Gaudete}, which gives the Sunday its name. Rose
vestments may be used today. Paul has more to ask of his readers than outward
cheerfulness. He tells them to bring their needs to God with thanksgiving
(Phil. 4:4--7). He can be grateful while there are still things he needs to
ask for.

In the Gospel the men sent from Jerusalem want John to say who he is. He calls
himself a voice crying in the wilderness. Then he tells them what they have
missed. “There hath stood one in the midst of you, whom you know not” (John
1:26). Gregory distinguishes the voice that is heard and falls silent from the
Word it makes known (Homily 7 on the Gospels, §§1--3). The people questioning
John can see Jesus as a man. They need faith to recognize that the one born of
Mary in time is the Son eternally begotten of the Father. John's own humility
lets him tell the truth about both himself and his Lord. “In that he arrogated
not to himself the name of Christ, he became a member of Christ.”

Gregory says that the way becomes straight in the heart of the person who
hears Christ's words with humility and puts them into practice. He quotes the
promise that the Father and the Son will make their dwelling with the one who
loves Christ and keeps his word (John 14:23). This helps explain why the
preparation for joy includes repentance. A person who can own a wrong and
leave it behind has begun to welcome the one John announces. At your evening
prayer this Sunday, thank God for some good you can name while you are still
asking for something you need.

\subsection{The winter Ember Days and the last evenings}

On Wednesday, Friday and Saturday after the Third Sunday, the Missal has
Masses for the winter Ember Days. The station on Wednesday is St Mary Major,
where the Gospel recounts Gabriel's visit to Mary (Luke 1:26--38). On Friday
at the Twelve Holy Apostles we hear how she went in haste to Elizabeth and how
the child leapt at her greeting (Luke 1:39--47). The fast of the season has
brought us to the beginning of the Lord's life among us. In both readings he
is present before anyone can see him.

The Saturday station is St Peter's. Among the readings of its longer Mass are
the three young men in the furnace and the angel of the Lord who makes the
middle of the fire like a wind bringing dew (Dan. 3:49--50). When we ask
heaven for dew with Isaiah on the Fourth Sunday, we can think of how God came
to the help of those young men where they were.

From December 17 to 23, whenever these dates fall in relation to the Ember
Days, the Church sings the seven O Antiphons at Vespers. These short chants
are the Advent antiphons for the Magnificat. Each addresses Christ by a name
drawn from Scripture and asks him to come. On December 20 the Church calls to
the Key of David, who can open what no one else can shut, and asks him to
bring out the prisoners who sit in darkness (Isa. 22:22, 42:7). Pray the
antiphon of the date, Mary's canticle and the antiphon again. As you say her
words about God's mercy to Abraham and his descendants, you are asking him to
keep the promise for which she is already giving thanks.

\subsection{The Fourth Sunday}

St Luke begins with the emperor Tiberius. He names Pilate and the other rulers
and the high priests. Then the word of God comes to John in the wilderness.
The voice is heard along the Jordan, so that people can go out to the river
and be baptized by him. John preaches a baptism of repentance, and Luke tells
what Isaiah said would happen when the Lord came. “Every valley shall be
filled and every mountain and hill shall be brought low” (Luke 3:1--6).

Gregory distinguishes John's baptism, which made ready for Christ, from
Christ's baptism, which forgives sins. When Gregory comes to the hills, he
pictures water running off high ground and springing up below. So the teaching
of God slips past the person raised up in pride and can be received by the one
who has become lowly. Gregory finds that humility in John, who refused to
claim Christ's place for himself (Homily 20 on the Gospels, §§2--4). Later in
Luke's account the people ask how to repent. John tells someone with two coats
to give to a person who has none. To a collector of taxes he says to take only
what is due (Luke 3:10--14). If you know what you owe someone, there is still
time before the feast to return it.

The station is the Twelve Holy Apostles. The Mass begins \emph{Rorate caeli},
“Drop down dew, ye heavens, from above” (Isa. 45:8). The earth is asked to
open and bring forth a saviour. In the Epistle Paul leaves the discovery of
the heart's secrets to the Lord, who will bring them to light when he comes (1
Cor. 4:1--5). At the offertory the words are those Gabriel spoke to Mary (Luke
1:28). In her the promised Saviour has been received and will come forth into
the light of day.
